%%%%%%%%%%%%%%%%%%%%%%%%%%%%%%%%%%%%%%%%%%%%%%%%%%%%%%%%%%%%%%%%%%%%%%
% Lớp 11 | Đường thẳng, mặt phẳng. Quan hệ song song trong | Hai mặt phẳng song song
% Dạng 1H4?4-7: Toán thực tế áp dụng hai mặt phẳng song song
% Số câu: 1  (sinh tự động bởi tach_cau.py; từ nay thầy sửa trực tiếp file này)
%%%%%%%%%%%%%%%%%%%%%%%%%%%%%%%%%%%%%%%%%%%%%%%%%%%%%%%%%%%%%%%%%%%%%%

%%%%%%%%%%%%%%%% TRẢ LỜI NGẮN %%%%%%%%%%%%%%%%

% nguồn: Ok_TraLoiNgan_Lop11_Lan2
\begin{ex}%[1H4H4-7]%[ma:eb0dccd6]
	Một kệ để đồ bằng gỗ có mâm tầng dưới $\left(ABCD\right)$ và mâm tầng trên $\left(EFGH\right)$  song song với nhau. Bác thợ mộc đo được $AE=100$ cm, $CG=120$ cm và muốn đóng thêm mâm tầng giữa $\left(IJKL\right)$ song song với hai mâm tầng trên, tầng dưới và $EI=42$ cm. Tính độ dài đoạn thẳng $KG$.
	\begin{center}
	 \definecolor{go}{RGB}{227,176,122}
		\begin{tikzpicture}[scale=1,font=\footnotesize,line join=round,line cap=round,>=stealth,declare function={k=0.7; l=0.6; h=3.5;}]
				 \path (0,0) coordinate (A)++(5,0) coordinate (D)++(1,1) coordinate (C)
					 ($(A)+(C)-(D)$) coordinate (B)
					($(A)!k!(D)$)  coordinate (a)
					($(B)!k!(C)$) coordinate  (b)
					($(A)+(0,h)$) coordinate  (E)
					($(B)+(0,h)$) coordinate (F) 
					($(b)+(0,h)$)  coordinate (G)
					($(a)+(0,h)$)  coordinate (H)
					($(A)!l!(E)$)  coordinate (I)
					 ($(D)!l!(H)$) coordinate (J)
					($(C)!l!(G)$)  coordinate (K)
					 ($(B)!l!(F)$) coordinate (L);
					 circle (1pt);
				\draw[line width=2pt] (B)--(F); 
				\draw[fill=go] (A)--(B)--(C)--(D)--cycle  (E)--(F)--(G)--(H)--cycle (I)--(J)--(K)--(L)--cycle;
				\draw[line width=2pt]  (A)--(E)
					 (D)--(H)
					 (C)--(G);
				\foreach \p/\g in {A/180,B/150,C/30,D/-45,E/180,F/90,G/60,H/0,I/180,J/0,K/0,L/45} \fill[black] (\p) circle (1pt)+(\g:0.3) node{$\p$}; 
			\end{tikzpicture}
	\end{center}
	\shortans{50,4} 
	\loigiai{ 
		Áp dụng  định lý Thales trong không gian, ta có $\dfrac{KG}{EI}=\dfrac{GC}{EA}\Leftrightarrow KG=\dfrac{42\cdot 120}{100}=50{,}4$.\\
	Vậy $KG=50{,}4$ cm.} 
\end{ex}
